% s03-1-5-approx.tex —— 3.1.5 小节：逼近（印刷页 77–81）

\subsection{逼近}
\label{sec:approximation}

本小节陈述并证明一个逼近结果，后文将用到它。该结果是 \cite[推论 4.5.7]{ABC+24a}
从因子化范畴到对（因子化范畴 $\mathcal{C}$ 与 $x_0$ 处的因子化 $\mathcal{C}$-模范畴
$\mathcal{M}$）的推广。

先重述该文 4.5 节的设定。

设 $Y$ 为 $X$ 上的仿射 $D$-概形，记 $T^+=L^+_\nabla Y$。若 $T$ 是 ind-placid
ind-概形的因子化空间，且
\[
\iota\colon T^+\to T
\]
是因子化空间的映射，并延拓为相对 $T^+$ 的 unital-in-correspondences 结构
（定义见 \cite[附录 C.10.6]{ABC+24a}），则
\[
\iota_*\mathcal{O}_{T^+}
\]
是因子化范畴 $\QCoh_{co}(T)$ 的因子化单位。

设 $Y_0$ 为 $D$-预叠，使得 $T^+_0=L^+_\nabla Y_0$ 具有仿射对角线。再设
\[
f\colon Y\to Y_0
\]
为 $X_{dR}$ 上的映射，使得
\[
L^+f\colon T^+=L^+_\nabla Y\to L^+_\nabla Y_0=T^+_0
\]
延拓为映射
\[
Lf\colon T\to T^+_0 .
\]
\srcpage{78}

有因子化范畴的宽松带单位函子
\[
\QCoh_{co}(T)\xrightarrow{(Lf)_*}\QCoh_{co}(T^+_0)\to \QCoh(T^+_0),
\]
它把因子化单位 $\iota_*\mathcal{O}_{T^+}$ 送到 $\QCoh(T^+_0)$ 中的因子化单位
$(L^+f)_*\mathcal{O}_{T^+}$。于是上述函子经带单位的因子化函子
\[
(Lf)^{\mathrm{enh}}_*\colon \QCoh_{co}(T)\to
(L^+f)_*\mathcal{O}_{T^+}\text{-}\FactMod\bigl(\QCoh(T^+_0)\bigr)
\]
分解。

下面回顾该文推论 4.5.7 的陈述。

\begin{proposition}[{\cite[推论 4.5.7]{ABC+24a}}]
\label{prop:approx-ABC}\label{prop:3.1.5.1}
设存在仿射 $D$-概形的映射
\[
\widetilde{Y}_0\to Y_0,
\]
它诱导 fpqc 覆盖
\[
\widetilde{T}^+_0=L^+_\nabla\widetilde{Y}_0\to L^+_\nabla Y_0=T^+_0 .
\]
再设存在等同
\[
L_\nabla(Y\times_{Y_0}\widetilde{Y}_0)\times_{L_\nabla(\widetilde{Y}_0)}
L^+_\nabla(\widetilde{Y}_0)
\simeq T\times_{T^+_0}\widetilde{T}^+_0,
\]
在此等同下，$\iota\colon T^+\to T$ 沿上述 fpqc 覆盖的基变换
\[
T^+\times_{T^+_0}\widetilde{T}^+_0\to T\times_{T^+_0}\widetilde{T}^+_0
\]
与映射
\[
T^+\times_{T^+_0}\widetilde{T}^+_0
=L^+_\nabla(Y\times_{Y_0}\widetilde{Y}_0)\to
L_\nabla(Y\times_{Y_0}\widetilde{Y}_0)\times_{L_\nabla(\widetilde{Y}_0)}
L^+_\nabla(\widetilde{Y}_0)
\]
相重合。最后设投影
\[
Y\times_{Y_0}\widetilde{Y}_0\to \widetilde{Y}_0
\]
是 D-afp 的。则因子化函子
\[
(Lf)^{\mathrm{enh}}_*\colon \QCoh_{co}(T)\to
(L^+f)_*\mathcal{O}_{T^+}\text{-}\FactMod\bigl(\QCoh(T^+_0)\bigr)
\]
在最终余连通子范畴上诱导等价
\[
\QCoh_{co}(T)^{>-\infty}
\simeq\bigl[(L^+f)_*\mathcal{O}_{T^+}\text{-}\FactMod
\bigl(\QCoh(T^+_0)\bigr)\bigr]^{>-\infty}.
\]
\end{proposition}
\srcpage{80}

下面详述我们对上述范式的改造。设 \cite[推论 4.5.7]{ABC+24a} 的假设有下述加强：
\[
T=L_\nabla Y\times_{L_\nabla Y_0}L^+_\nabla Y_0 .
\]
设 $S_0$ 为带映射
\[
S_0\to (L_\nabla Y_0)_{\Ran_{x_0}}
\]
的因子化 $T_0$-模空间。则
\[
(L_\nabla Y)_{\Ran_{x_0}}\times_{(L_\nabla Y_0)_{\Ran_{x_0}}}S_0
\]
是因子化 $T$-模空间，并有对的态射
\begin{align*}
\Bigl(&\QCoh_{co}\bigl(L_\nabla Y\times_{L_\nabla Y_0}L^+_\nabla Y_0\bigr),\\
&\QCoh_{co}\bigl((L_\nabla Y)_{\Ran_{x_0}}
\times_{(L_\nabla Y_0)_{\Ran_{x_0}}}S_0\bigr)\Bigr)
\to \bigl(\QCoh(L^+_\nabla Y_0),\QCoh(S_0)\bigr),
\end{align*}
它提升为
\begin{equation}
\begin{aligned}
\Bigl(&\QCoh_{co}\bigl(L_\nabla Y\times_{L_\nabla Y_0}L^+_\nabla Y_0\bigr),\\
&\QCoh_{co}\bigl((L_\nabla Y)_{\Ran_{x_0}}
\times_{(L_\nabla Y_0)_{\Ran_{x_0}}}S_0\bigr)\Bigr)\\
&\qquad\to
(L^+f)_*\mathcal{O}_{T^+}\text{-}\FactMod
\bigl(\QCoh(L^+_\nabla Y_0),\QCoh(S_0)\bigr).
\end{aligned}
\label{eq:3.1.1}
\end{equation}

我们断言

\begin{theorem}
\label{thm:3.1.5.2}
设 $Y_0=\pt/H\times X_{dR}$，$H$ 为仿射代数群。则式~\eqref{eq:3.1.1} 中的态射
在每一分量的最终余连通子范畴上诱导等价。
\end{theorem}

\begin{proof}
$H=\{e\}$ 的情形显然。一般情形取 $\widetilde{Y}_0=\pt\times X_{dR}$，并记
\[
\widetilde{S}_0=(L_\nabla\widetilde{Y}_0)_{\Ran_{x_0}}
\times_{(L_\nabla Y_0)_{\Ran_{x_0}}}S_0 .
\]
再记 $\widetilde{S}_0\to S_0$ 的 Čech 神经为 $\widetilde{S}^{\bullet}_0$。按该文的记号有
\[
\QCoh_{co}\Bigl((L_\nabla Y)_{\Ran_{x_0}}
\times_{(L_\nabla Y_0)_{\Ran_{x_0}}}S_0\Bigr)
\simeq \Tot\Bigl(\QCoh_{co}\bigl((L_\nabla\widetilde{Y}^{\bullet})_{\Ran_{x_0}}
\times_{(L_\nabla\widetilde{Y}^{\bullet}_0)_{\Ran_{x_0}}}\widetilde{S}^{\bullet}_0\bigr)\Bigr)
\]
以及
\[
(L^+f)_*\mathcal{O}_{T^+}\text{-}\FactMod\bigl(\QCoh(S_0)\bigr)
\simeq \Tot\Bigl((L^+\tilde f^{\bullet})_*\mathcal{O}_{\widetilde{T}^{+,\bullet}}
\text{-}\FactMod\bigl(\QCoh(\widetilde{S}^{\bullet}_0)\bigr)\Bigr).
\]
因此结论由 $H=\{e\}$ 的情形（及其拉回）得出。
\end{proof}
\srcpage{81}
